% Part: normal-modal-logic
% Chapter: filtrations
% Section: more-filtrations

\documentclass[../../../include/open-logic-section]{subfiles}

\begin{document}

\olfileid{nml}{fil}{acc}

\olsection{Filtrasi lan Sipat Aksesibilitas}

Kaya kang wis kasebut, filtrasi saka model universal, serial, lan
refleksif mesthi uga universal, serial, utawa refleksif. Nanging ora
saben filtrasi saka model simetris utawa transitif uga simetris utawa
transitif, miturut urutane. Ing sawetara kasus, filtrasi bisa
ditepakake supaya sipat iku tetep kawujud. Carane padha karo definisi
filtrasi paling kasar, nanging kita nambah sarat ing definisi~$R^*$.
Upamane $\Gamma$ katutup tumrap sub-!!{formula}s. Gatekna
relasi~$C_i(u,v)$ ing \olref{tab:Cn-filtrations} antarane jagad
$u$, $v$ ing model $\mModel{M} =\tuple{W, R, V}$. Kita bisa netepake
$R^*[u][v]$ adhedhasar gabungan sarat-sarat iki. Upamane, yen kita
netepake $R^*[u][v]$ yen lan mung yen sarat $C_1(u,v)$ kawujud,
kita entuk persis filtrasi paling kasar. Yen kita netepake
$R^*[u][v]$ yen lan mung yen $C_1(u,v)$ lan $C_2(u, v)$ loro-lorone
kawujud, kita entuk filtrasi liyane. Filtrasi iki ``luwih alus''
tinimbang filtrasi paling kasar amarga luwih sithik pasangan jagad
kang nyukupi $C_1(u,v)$ lan $C_2(u,v)$ tinimbang kang mung nyukupi
$C_1(u,v)$.

\begin{table}[ht]
  \centering
  \begin{tabular}{|ll|}
    \hline
    \multirow{2}{*}{$C_1(u,v)$:}
    \iftag{prvBox}{&
      yen $\Box!A \in \Gamma$ lan $\mSat{M}{\Box!A}[u]$,
      mula $\mSat{M}{!A}[v]$;\iftag{prvDiamond}{ lan}{}\\}{}
    \iftag{prvDiamond}{&
      yen $\Diamond!A \in \Gamma$ lan $\mSat{M}{!A}[v]$,
      mula $\mSat{M}{\Diamond!A}[u]$; \\}{}
    \hline
    \multirow{2}{*}{$C_2(u,v)$:} 
    \iftag{prvBox}{&
      yen $\Box!A \in \Gamma$ lan $\mSat{M}{\Box!A}[v]$,
      mula $\mSat{M}{!A}[u]$;\iftag{prvDiamond}{ lan}{}\\}{}
    \iftag{prvDiamond}{&
      yen $\Diamond!A \in \Gamma$ lan $\mSat{M}{!A}[u]$,
      mula $\mSat{M}{\Diamond!A}[v]$; \\}{}
    \hline
    \multirow{2}{*}{$C_3(u,v)$:}
    \iftag{prvBox}{&
      yen $\Box!A \in \Gamma$ lan $\mSat{M}{\Box!A}[u]$,
      mula $\mSat{M}{\Box!A}[v]$;\iftag{prvDiamond}{ lan}{}\\}{}
    \iftag{prvDiamond}{&
      yen $\Diamond!A \in \Gamma$ lan $\mSat{M}{\Diamond!A}[v]$,
      mula $\mSat{M}{\Diamond!A}[u]$; \\}{}
    \hline
    \multirow{2}{*}{$C_4(u,v)$:}
    \iftag{prvBox}{&
      yen $\Box!A \in \Gamma$ lan $\mSat{M}{\Box!A}[v]$,
      mula $\mSat{M}{\Box!A}[u]$;\iftag{prvDiamond}{ lan}{}\\}{}
    \iftag{prvDiamond}{&
      yen $\Diamond!A \in \Gamma$ lan $\mSat{M}{\Diamond!A}[u]$,
      mula $\mSat{M}{\Diamond!A}[v]$; \\}{}
    \hline
    \end{tabular}
  \caption{Sarat ing jagad-jagad kang mungkin kanggo netepake
    filtrasi.}
  \ollabel{tab:Cn-filtrations}
\end{table}

\begin{thm}\ollabel{thm:more-filtrations}
  Upamane $\mModel{M} =\tuple{W,R,V}$ sawijining model, lan $\Gamma$
  katutup tumrap sub-!!{formula}s. Tetepna $W^*$ lan $V^*$ kaya ing
  \olref[fil]{defn:filtration}. Banjur:
  \begin{enumerate}
  \item Upamane $R^*[u][v]$ yen lan mung yen $C_1(u, v) \land
    C_2(u,v)$. Banjur $R^*$ simetris, lan $\mModel{M^*} =
    \tuple{W^*,R^*,V^*}$ iku filtrasi yen $\mModel{M}$ simetris.
  \item Upamane $R^*[u][v]$ yen lan mung yen $C_1(u, v)
    \land C_3(u,v)$. Banjur $R^*$ transitif, lan
    $\mModel{M^*}=\tuple{W^*,R^*,V^*}$ iku filtrasi yen
    $\mModel{M}$ transitif.
  \item Upamane $R^*[u][v]$ yen lan mung yen $C_1(u, v) \land C_2(u,v)
    \land C_3(u,v) \land C_4(u,v)$. Banjur $R^*$ simetris lan
    transitif, lan $\mModel{M^*}=\tuple{W^*,R^*,V^*}$ iku filtrasi
    yen $\mModel{M}$ simetris lan transitif.
  \item Upamane $R^*$ ditepakake minangka $R^*[u][v]$ yen lan mung yen
    $C_1(u, v) \land C_3(u,v) \land C_4(u,v)$. Banjur $R^*$ transitif
    lan Euklidean, lan $\mModel{M^*}=\tuple{W^*,R^*,V^*}$ iku filtrasi
    yen $\mModel{M}$ transitif lan Euklidean.
  \end{enumerate}
\end{thm}

\begin{proof}
  \begin{enumerate}
    \item Langsung cetha yen $R^*$ simetris, amarga $C_1(u,v)
      \Leftrightarrow C_2(v,u)$ lan $C_2(u,v) \Leftrightarrow
      C_1(v,u)$. Dadi isih kudu dituduhake yen $\mModel{M}$ simetris,
      mula $\mModel{M^*}$ iku filtrasi liwat~$\Gamma$. Sarat
      $C_1(u,v)$ njamin yen
      \iftag{prvBox}{%
        \iftag{prvDiamond}{%
          \olref[fil]{defn:filtration-R2} lan
          \olref[fil]{defn:filtration-R3} saka
          \olref[fil]{defn:filtration} padha}{%
          \olref[fil]{defn:filtration}\olref[fil]{defn:filtration-R2}}}{%
        \olref[fil]{defn:filtration}\olref[fil]{defn:filtration-R3}}
      kawujud. Dadi kita mung kudu mriksa
      \olref[fil]{defn:filtration}\olref[fil]{defn:filtration-R1},
      yaiku yen $Ruv$ njalari $R^*[u][v]$.

      Dadi upamane $Ruv$. Kanggo nuduhake $R^*[u][v]$, kita kudu
      netepake $C_1(u,v)$ lan $C_2(u,v)$. Kanggo $C_1$:
      \iftag{prvBox}{yen $\Box!A \in\Gamma$ lan
        $\mSat{M}{\Box!A}[u]$, mula uga $\mSat{M}{!A}[v]$ (amarga
        $Ruv$).\iftag{prvDiamond}{ Kajaba iku,
        }{}}{}\iftag{prvDiamond}{yen $\Diamond!A \in \Gamma$ lan
        $\mSat{M}{!A}[v]$, mula $\mSat{M}{\Diamond!A}[u]$ amarga
        $Ruv$.}{} Kanggo $C_2$: \iftag{prvBox}{yen $\Box!A \in \Gamma$
        lan $\mSat{M}{\Box!A}[v]$, mula $Ruv$ njalari $Rvu$ miturut
        simetri, saengga $\mSat{M}{!A}[u]$.\iftag{prvDiamond}{ Kajaba
        iku, }{}}{}\iftag{prvDiamond}{yen $\Diamond!A \in\Gamma$ lan
        $\mSat{M}{!A}[u]$, mula $\mSat{M}{\Diamond!A}[v]$ (amarga
        $Rvu$ miturut simetri).}{}
    \item Latihan.
    \item Latihan.
    \item Latihan.
  \end{enumerate}
\end{proof}

\begin{prob}
  Jangkepna bukti \olref[nml][fil][acc]{thm:more-filtrations}.
\end{prob}

\end{document}
